\documentclass[UTF8,zihao=-4]{ctexart}
\usepackage[a4paper,margin=2.2cm]{geometry}
\usepackage{fontspec}
\usepackage{tabularx}
\usepackage{array}
\usepackage{ragged2e}
\usepackage{fancyhdr}
\usepackage{hyperref}
\setCJKmainfont{Noto Serif SC}
\setCJKsansfont{Noto Sans SC}
\setmainfont{Noto Serif SC}
\setlength{\parindent}{0pt}
\setlength{\parskip}{0.45em}
\pagestyle{fancy}
\fancyhf{}
\fancyhead[L]{商务智能}
\fancyfoot[C]{\thepage}
\hypersetup{colorlinks=true,linkcolor=black,urlcolor=blue}
\begin{document}
\title{3、ETL的作用}
\author{}
\date{}
\maketitle

ETL是抽取、转换、装载\par

1. 辨识与主题相关的原始数据\par
也就是从多个数据源中找出对分析主题有用的数据。\par

2. 开发数据抽取策略\par
解决从哪里抽、抽哪些、什么时候抽、怎么抽的问题。\par

3. 保证数据正确和完整\par
数据进入仓库前，需要处理错误、缺失、不一致等问题。\par

4. 将原始数据转换为目标规格\par
包括格式转换、编码转换、单位统一、字段重命名、关键字转换、数据清洗、数据合并、数据汇总等。\par

5. 将数据加载到预定目标区域\par
目标可以包括原子层和集成数据、数据集市、ODS、缓冲区等。\par

ETL 的作用是把来自不同数据源的原始数据抽取出来，经过清洗、转换、集成、汇总等处理，使其变成正确、完整、一致、符合目标规格的数据，并装载到数据仓库、数据集市、ODS 或缓冲区中，为后续 OLAP、报表查询和数据挖掘提供可靠的数据基础。\par

\end{document}
